% ============================================================================
%  preamble_lecture_notes.tex
%  A ready-to-use, professional lecture-notes preamble for the `report` class.
%
%  Distilled and GENERALIZED from a real, tested preamble that compiled a
%  131-page lecture book cleanly (a terse PRX paper expanded into self-contained
%  notes). Everything here is PROJECT-AGNOSTIC; the only thing you customise is
%  the clearly-marked block near the bottom (title-page subtitle + domain math
%  macros). Drop this file next to your master and `\input` it.
%
%  --------------------------------------------------------------------------
%  THE KEY TRICK  (read this once -- it is why this preamble is reusable)
%  --------------------------------------------------------------------------
%  Theorems are declared with plain `amsthm` (\newtheorem); the colored, tinted,
%  left-barred BOX around each is added PURELY as a visual layer, by wrapping
%  the existing amsthm environment with tcolorbox's \tcolorboxenvironment. The
%  environment NAMES, COUNTERS, and \label scheme are unchanged, so the calling
%  convention in your chapters stays the ORDINARY one:
%
%        \begin{theorem}\label{thm:foo} ... \end{theorem}
%
%    * Do NOT convert theorems to \newtcbtheorem / \newtcolorbox theorem boxes:
%      that changes the calling convention (\begin{theorem}{title}{label}) and
%      breaks every chapter. Keep amsthm + \tcolorboxenvironment.
%    * If a compile error arises, fix the offending chapter construct, not this
%      design. Do not "simplify" the preamble away during assembly.
%
%  Full WHY -- the layering discipline, package stack and load order, font and
%  heading design, the sanctioned 5-color palette, and failure modes -- lives in
%  references/typesetting_guide.md, the source of truth for this design.
%  (Palette hex values are defined once, in the COLOR PALETTE block below.)
%
%  Requires a reasonably complete TeX Live (newpx, tgheros, tcolorbox[most],
%  titlesec, microtype, beramono). To locate pdflatex, probe in order:
%  `command -v pdflatex`; else /Library/TeX/texbin/pdflatex (macOS MacTeX);
%  else /usr/local/texlive/*/bin/*/pdflatex (Linux TeX Live).
% ============================================================================

% ----- encoding, fonts, typography ------------------------------------------
\usepackage[T1]{fontenc}
\usepackage[utf8]{inputenc}
\usepackage{amsmath}
\usepackage{amssymb}
\usepackage{mathtools}
\usepackage{amsthm}
\usepackage{tgheros}                 % Helvetica-like sans: headings, captions, labels
\usepackage{newpxtext}               % Palatino-like serif body text
\usepackage{newpxmath}               % matching math (load AFTER newpxtext)
\usepackage{bm}
\usepackage[scaled=0.86]{beramono}   % clean mono for code/paths (optional use)
\usepackage{microtype}               % refined spacing/kerning -- load late
\linespread{1.06}                    % a touch more leading for readability

% ----- page geometry ---------------------------------------------------------
\usepackage[a4paper,top=1.15in,bottom=1.2in,inner=1.25in,outer=1.25in,
            headheight=14pt,marginparwidth=0.9in]{geometry}

% ----- graphics, color, tables, lists ---------------------------------------
\usepackage{graphicx}
\usepackage{xcolor}
\usepackage{booktabs}
\usepackage{colortbl}          % row/cell color for styled results tables (archetype I)
\usepackage{enumitem}
\usepackage{caption}
\usepackage[most]{tcolorbox}
\usepackage{tikz}
\usetikzlibrary{arrows.meta,calc,positioning,decorations.pathmorphing,%
  backgrounds,fit,patterns,shapes.geometric,plotmarks}
\usepackage{titlesec}
\usepackage{fancyhdr}
\usepackage{needspace}

% hyperref late; bookmarks; soft colored links
\usepackage[colorlinks=true,linktoc=all]{hyperref}

% ----------------------------------------------------------------------------
%  COLOR PALETTE  (the only sanctioned colors)
% ----------------------------------------------------------------------------
\definecolor{navy}{HTML}{1f3b73}   % primary   / structure / theorems
\definecolor{mpred}{HTML}{c0392b}  % emphasis  / "hot" side / warnings
\definecolor{teal}{HTML}{16887b}   % secondary / definitions / edges
\definecolor{gold}{HTML}{b8860b}   % tertiary  / examples / highlights / zeros
\definecolor{inkgray}{HTML}{4a4a4a} % neutral / running headers / subsubsections / remark & proof-sketch rules

\hypersetup{
  linkcolor=navy,
  citecolor=teal,
  urlcolor=mpred,
  filecolor=navy,
  bookmarksnumbered=true,
}

% ----------------------------------------------------------------------------
%  HEADINGS  (titlesec) -- modern colored chapter & section design
%  Chapter opener: big faded navy chapter number, right-aligned, over a rule.
% ----------------------------------------------------------------------------
\titleformat{\chapter}[display]
  {\sffamily\bfseries\color{navy}}
  {\raggedleft\sffamily\bfseries\fontsize{56}{56}\selectfont\color{navy!22}%
     \thechapter}
  {6pt}
  {\Huge\raggedright}
  [\vspace{2pt}{\color{navy!22}\titlerule[1.2pt]}]
\titlespacing*{\chapter}{0pt}{6pt}{22pt}

\titleformat{\section}
  {\sffamily\Large\bfseries\color{navy}}{\thesection}{0.8em}{}
\titleformat{\subsection}
  {\sffamily\large\bfseries\color{navy!88}}{\thesubsection}{0.7em}{}
\titleformat{\subsubsection}
  {\sffamily\normalsize\bfseries\color{inkgray}}{\thesubsubsection}{0.6em}{}
\titlespacing*{\section}{0pt}{14pt}{6pt}
\titlespacing*{\subsection}{0pt}{11pt}{4pt}

% ----------------------------------------------------------------------------
%  RUNNING HEADERS / FOOTERS  (fancyhdr, oneside)
% ----------------------------------------------------------------------------
\pagestyle{fancy}
\fancyhf{}
\renewcommand{\headrulewidth}{0.4pt}
\renewcommand{\footrulewidth}{0pt}
\renewcommand{\chaptermark}[1]{\markboth{\thechapter.\ #1}{}}
\renewcommand{\sectionmark}[1]{\markright{\thesection\ #1}}
\fancyhead[L]{\sffamily\footnotesize\color{inkgray}\nouppercase{\leftmark}}
\fancyhead[R]{\sffamily\footnotesize\color{inkgray}\nouppercase{\rightmark}}
\fancyfoot[C]{\sffamily\small\color{navy}\thepage}
% chapter-opening pages use the `plain` style
\fancypagestyle{plain}{%
  \fancyhf{}\renewcommand{\headrulewidth}{0pt}%
  \fancyfoot[C]{\sffamily\small\color{navy}\thepage}}

% ----------------------------------------------------------------------------
%  CAPTIONS  (sans, bold colored label -- Nature-ish)
% ----------------------------------------------------------------------------
\captionsetup{font={small},labelfont={bf,sf,color=navy},
  textfont={sf},labelsep=period,width=0.95\linewidth}

% ----------------------------------------------------------------------------
%  THEOREM ENVIRONMENTS  (amsthm names/counters preserved; tcolorbox styling)
%  Step 1: colored bold sans amsthm header styles.
%  Step 2: declare the environments with \newtheorem (ordinary calling convention).
%  Step 3: add the visual box with \tcolorboxenvironment (THE KEY TRICK).
% ----------------------------------------------------------------------------
% --- Step 1: custom amsthm styles (colored bold headers) ---
\newtheoremstyle{thmnavy}{8pt}{6pt}{\itshape}{0pt}
  {\sffamily\bfseries\color{navy}}{.}{0.6em}{}
\newtheoremstyle{thmteal}{8pt}{6pt}{\normalfont}{0pt}
  {\sffamily\bfseries\color{teal!85!black}}{.}{0.6em}{}
\newtheoremstyle{thmgold}{8pt}{6pt}{\normalfont}{0pt}
  {\sffamily\bfseries\color{gold!72!black}}{.}{0.6em}{}
\newtheoremstyle{thmremark}{8pt}{6pt}{\normalfont}{0pt}
  {\itshape\color{inkgray}}{.}{0.5em}{}

% --- Step 2: declare environments (shared counter, numbered within chapter) ---
\theoremstyle{thmnavy}
\newtheorem{theorem}{Theorem}[chapter]
\newtheorem{proposition}[theorem]{Proposition}
\newtheorem{lemma}[theorem]{Lemma}
\newtheorem{corollary}[theorem]{Corollary}

\theoremstyle{thmteal}
\newtheorem{definition}[theorem]{Definition}

\theoremstyle{thmgold}
\newtheorem{example}[theorem]{Example}

\theoremstyle{thmremark}
\newtheorem{remark}[theorem]{Remark}

% convenience alias: \begin{defn}...\end{defn}
\newenvironment{defn}{\begin{definition}}{\end{definition}}

% --- Step 3: visual boxes -- wrap the amsthm environments (THE KEY TRICK).
% navy = theorem/proposition/lemma/corollary; teal = definition;
% gold = example; gray = remark. Calling convention stays \begin{...}\end{...}.
\tcolorboxenvironment{theorem}{
  enhanced jigsaw, breakable, before skip=10pt, after skip=10pt,
  boxrule=0pt, frame hidden, sharp corners,
  colback=navy!5, borderline west={3pt}{0pt}{navy},
  left=12pt,right=10pt,top=7pt,bottom=7pt}
\tcolorboxenvironment{proposition}{
  enhanced jigsaw, breakable, before skip=10pt, after skip=10pt,
  boxrule=0pt, frame hidden, sharp corners,
  colback=navy!5, borderline west={3pt}{0pt}{navy},
  left=12pt,right=10pt,top=7pt,bottom=7pt}
\tcolorboxenvironment{lemma}{
  enhanced jigsaw, breakable, before skip=10pt, after skip=10pt,
  boxrule=0pt, frame hidden, sharp corners,
  colback=navy!5, borderline west={3pt}{0pt}{navy},
  left=12pt,right=10pt,top=7pt,bottom=7pt}
\tcolorboxenvironment{corollary}{
  enhanced jigsaw, breakable, before skip=10pt, after skip=10pt,
  boxrule=0pt, frame hidden, sharp corners,
  colback=navy!5, borderline west={3pt}{0pt}{navy},
  left=12pt,right=10pt,top=7pt,bottom=7pt}
\tcolorboxenvironment{definition}{
  enhanced jigsaw, breakable, before skip=10pt, after skip=10pt,
  boxrule=0pt, frame hidden, sharp corners,
  colback=teal!6, borderline west={3pt}{0pt}{teal!85!black},
  left=12pt,right=10pt,top=7pt,bottom=7pt}
\tcolorboxenvironment{example}{
  enhanced jigsaw, breakable, before skip=10pt, after skip=10pt,
  boxrule=0pt, frame hidden, sharp corners,
  colback=gold!7, borderline west={3pt}{0pt}{gold!72!black},
  left=12pt,right=10pt,top=7pt,bottom=7pt}
\tcolorboxenvironment{remark}{
  enhanced jigsaw, breakable, before skip=8pt, after skip=8pt,
  boxrule=0pt, frame hidden, sharp corners,
  colback=black!3, borderline west={2.2pt}{0pt}{inkgray},
  left=11pt,right=8pt,top=5pt,bottom=5pt}

% ----------------------------------------------------------------------------
%  BOXED-RESULT ENVIRONMENT  (keyresult tcolorbox -- headline result of a chapter)
%  Use sparingly: at most one or two per chapter.
% ----------------------------------------------------------------------------
\newtcolorbox{keyresult}[1][]{%
  enhanced, breakable,
  colback=gold!7,
  colframe=gold!75!black,
  boxrule=0.9pt,
  arc=2.5pt,
  left=8pt,right=8pt,top=8pt,bottom=8pt,
  fonttitle=\sffamily\bfseries\color{white},
  coltitle=white,
  attach boxed title to top left={yshift=-2.6mm,xshift=5mm},
  boxed title style={colback=gold!75!black,arc=1pt},
  title={Key result},
  #1,
}

% ----------------------------------------------------------------------------
%  GENERIC MATH MACROS  (domain-agnostic -- safe to keep on any project)
% ----------------------------------------------------------------------------
\newcommand{\bL}{\mathbf{L}}              % a bold matrix L (Laplacian/operator)
\newcommand{\bI}{\mathbf{I}}              % identity
\newcommand{\bone}{\mathbf{1}}            % all-ones vector
\newcommand{\adj}{\operatorname{adj}}     % adjugate
\newcommand{\Cov}{\operatorname{Cov}}     % covariance
\newcommand{\sgn}{\operatorname{sgn}}     % sign
\newcommand{\diag}{\operatorname{diag}}   % diagonal matrix
\newcommand{\Tr}{\operatorname{Tr}}       % trace
\newcommand{\Res}{\operatorname{Res}}     % residue
\providecommand{\R}{}\renewcommand{\R}{\mathbb{R}} % reals (robust if predefined)
\providecommand{\C}{}\renewcommand{\C}{\mathbb{C}} % complex  (robust if predefined)
\providecommand{\N}{}\renewcommand{\N}{\mathbb{N}} % naturals (robust if predefined)
\providecommand{\Z}{}\renewcommand{\Z}{\mathbb{Z}} % integers (robust if predefined)
\newcommand{\ee}{\mathrm{e}}              % Euler's number
\newcommand{\dd}{\mathrm{d}}              % differential d
\newcommand{\eps}{\varepsilon}
\newcommand{\la}{\lambda}
\DeclarePairedDelimiter{\abs}{\lvert}{\rvert}
\DeclarePairedDelimiter{\norm}{\lVert}{\rVert}

% ============================================================================
%  ====  CUSTOMISE PER PROJECT  ===============================================
%  Everything ABOVE is project-agnostic. Edit ONLY this block per document.
% ============================================================================

% --- (a) Title-page subtitle: the italic line under the title rule.
%     Set this in your master with \renewcommand{\notessubtitle}{...} BEFORE
%     \maketitle, or edit the default here.
\providecommand{\notessubtitle}{A self-contained set of lecture notes}

% --- (b) Domain-specific math macros for YOUR paper. Examples (delete/replace):
%   \newcommand{\bpi}{\bm{\pi}}              % a bold Greek vector
%   \newcommand{\Psimat}{\bm{\Psi}}          % a named bold matrix
%   \newcommand{\Zw}{\mathcal{Z}_\omega}     % a named scalar / partition function
%   \newcommand{\osc}{\operatorname{osc}}    % a domain operator
% (Define them here so chapters can use a single fixed notation.)

% ============================================================================
%  ====  END CUSTOMISE PER PROJECT  ===========================================
% ============================================================================

% ----------------------------------------------------------------------------
%  TITLE PAGE  (custom; respects \title/\author/\date set in the master)
% ----------------------------------------------------------------------------
\makeatletter
\renewcommand{\maketitle}{%
  % titlepage resets the page counter; suppress its hyperref page anchor so it
  % does not collide with the real page 1 (avoids a pdfTeX duplicate-destination warning).
  \hypersetup{pageanchor=false}%
  \begin{titlepage}
  \centering
  \vspace*{0.16\textheight}
  {\color{navy}\rule{\textwidth}{1.4pt}}\\[1.4em]
  {\sffamily\bfseries\color{navy}\fontsize{30}{34}\selectfont \@title \par}
  \vspace{1.1em}
  {\color{navy}\rule{\textwidth}{1.4pt}}\\[2.4em]
  {\large\itshape \notessubtitle \par}
  \vfill
  {\sffamily\large\@author \par}
  \vspace{0.7em}
  {\large \@date \par}
  \vspace*{0.10\textheight}
  \end{titlepage}%
  \hypersetup{pageanchor=true}}
\makeatother

% ----------------------------------------------------------------------------
%  misc
% ----------------------------------------------------------------------------
\graphicspath{{figs/}{../figs/}{./}}
\setlength{\parindent}{0pt}
\setlength{\parskip}{5pt plus 1pt}
% prefer slightly loose lines over margin-breaking overfull boxes
% (TeX's emergency pass only runs when normal line breaking fails)
\setlength{\emergencystretch}{1.5em}
\setlist{itemsep=2pt,topsep=4pt}
\allowdisplaybreaks
